\documentclass[11pt]{article}
\usepackage[margin=1in]{geometry}
\usepackage{amsmath}
\usepackage{booktabs}
\usepackage{array}
\usepackage{xcolor}
\usepackage{hyperref}
\hypersetup{colorlinks=true, linkcolor=blue, urlcolor=blue}

\title{Normalized Source-Shape Fit Functions for EPOS Pion Pair Sources}
\author{}
\date{June 1, 2026}

\begin{document}
\maketitle

\section*{Purpose}

This note summarizes the normalized fit functions implemented in
\texttt{marco\_fitEPOSPrimordialLevy.C} for describing the EPOS primordial
pion pair source distributions. The goal is to compare different
non-Gaussian source-shape parameterizations without allowing an arbitrary
overall amplitude or pedestal term to improve the fit artificially.

\section*{Common Normalization Convention}

The input source histogram is normalized as a probability density:
\begin{equation}
  \int h(x)\,dx = 1.
\end{equation}

Each fit function is also normalized over the fit range:
\begin{equation}
  S_{\rm fit}(x;\vec p)
  =
  \frac{S_{\rm raw}(x;\vec p)}
       {\int_{x_{\rm min}}^{x_{\rm max}} S_{\rm raw}(x';\vec p)\,dx'}.
\end{equation}

For signed one-dimensional projections, such as
\(r_{\rm out}\), \(r_{\rm side}\), and \(r_{\rm long}\), the raw function is
used directly. For radial distributions, such as \(r_{\rm inv}^{*}\) or
\(|r|_{\rm LCMS}\), the raw function includes the radial shell factor:
\begin{equation}
  S_{\rm raw}^{\rm radial}(r)
  =
  r^2\,S_{\rm shape}(r).
\end{equation}

The factor \(4\pi\) is absorbed by normalization and therefore does not affect
the fitted shape parameters.

\section*{Fit Functions}

\subsection*{1. Normalized Double Gaussian}

For a signed 1D projection:
\begin{equation}
  S_{\rm raw}(x)
  =
  \exp\!\left(-\frac{x^2}{2R_{\rm narrow}^2}\right)
  +
  w_{\rm broad}
  \exp\!\left(-\frac{x^2}{2R_{\rm broad}^2}\right).
\end{equation}

For a radial source distribution:
\begin{equation}
  S_{\rm raw}(r)
  =
  r^2
  \left[
  \exp\!\left(-\frac{r^2}{2R_{\rm narrow}^2}\right)
  +
  w_{\rm broad}
  \exp\!\left(-\frac{r^2}{2R_{\rm broad}^2}\right)
  \right].
\end{equation}

\begin{center}
\begin{tabular}{lll}
\toprule
Parameter & Meaning & Comment \\
\midrule
\(R_{\rm narrow}\) & Width of the narrow Gaussian part & Describes the central peak scale \\
\(R_{\rm broad}\) & Width of the broad Gaussian part & Captures a broad tail-like shape \\
\(w_{\rm broad}\) & Relative broad Gaussian weight & Shape weight, not a physical source fraction \\
\bottomrule
\end{tabular}
\end{center}

This function is useful as a compact empirical comparison. However, the broad
Gaussian should not be interpreted as a separate pion source component. It is
only a basis term for describing a non-Gaussian averaged source shape.

\subsection*{2. Normalized Symmetric Levy}

For a signed 1D projection:
\begin{equation}
  S_{\rm raw}(x)
  =
  \exp\!\left[-\left(\frac{|x|}{R}\right)^\alpha\right].
\end{equation}

For a radial source distribution:
\begin{equation}
  S_{\rm raw}(r)
  =
  r^2
  \exp\!\left[-\left(\frac{r}{R}\right)^\alpha\right].
\end{equation}

\begin{center}
\begin{tabular}{lll}
\toprule
Parameter & Meaning & Comment \\
\midrule
\(R\) & Characteristic source scale & Homogeneity-length scale \\
\(\alpha\) & Levy stability / tail parameter & \(\alpha=2\) Gaussian, \(\alpha=1\) exponential-like \\
\bottomrule
\end{tabular}
\end{center}

This is the simplest normalized non-Gaussian model. In the current EPOS source
studies, it is often not flexible enough to describe both the sharp central
peak and the far tail simultaneously.

\subsection*{3. Normalized Gaussian Plus Exponential Tail}

For a signed 1D projection:
\begin{equation}
  S_{\rm raw}(x)
  =
  \exp\!\left(-\frac{x^2}{2R_G^2}\right)
  +
  w_{\rm exp}
  \exp\!\left(-\frac{|x|}{L_{\rm exp}}\right).
\end{equation}

For a radial source distribution:
\begin{equation}
  S_{\rm raw}(r)
  =
  r^2
  \left[
  \exp\!\left(-\frac{r^2}{2R_G^2}\right)
  +
  w_{\rm exp}
  \exp\!\left(-\frac{r}{L_{\rm exp}}\right)
  \right].
\end{equation}

\begin{center}
\begin{tabular}{lll}
\toprule
Parameter & Meaning & Comment \\
\midrule
\(R_G\) & Gaussian core width & Central source scale \\
\(L_{\rm exp}\) & Exponential tail length & Effective long-range falloff scale \\
\(w_{\rm exp}\) & Relative exponential tail weight & Shape weight, not a source fraction \\
\bottomrule
\end{tabular}
\end{center}

This function is motivated by the observation that the intermediate-to-large
distance region can appear approximately linear on a log-scale plot. The
exponential term should be viewed as an effective tail shape, not as a new
physical pion component by itself.

\subsection*{4. Normalized Log-Normal Radius Mixture}

The source is modeled as a continuous superposition of Gaussian sources with a
fluctuating radius:
\begin{equation}
  S_{\rm raw}(x)
  =
  \int dR\,
  P_{\rm LN}(R;R_{\rm med},\sigma_{\ln R})
  \exp\!\left(-\frac{x^2}{2R^2}\right),
\end{equation}
where
\begin{equation}
  P_{\rm LN}(R;R_{\rm med},\sigma_{\ln R})
  =
  \frac{1}{R\,\sigma_{\ln R}\sqrt{2\pi}}
  \exp\!\left[
    -\frac{\left(\ln R-\ln R_{\rm med}\right)^2}
          {2\sigma_{\ln R}^2}
  \right].
\end{equation}

For radial distributions, the same shell factor is applied:
\begin{equation}
  S_{\rm raw}(r)
  =
  r^2
  \int dR\,
  P_{\rm LN}(R;R_{\rm med},\sigma_{\ln R})
  \exp\!\left(-\frac{r^2}{2R^2}\right).
\end{equation}

\begin{center}
\begin{tabular}{lll}
\toprule
Parameter & Meaning & Comment \\
\midrule
\(R_{\rm med}\) & Median Gaussian radius & Typical homogeneity length \\
\(\sigma_{\ln R}\) & Width of radius fluctuations in log-space & Controls non-Gaussian tail strength \\
\bottomrule
\end{tabular}
\end{center}

The physical motivation is that the observed source is an event-averaged
superposition of approximately Gaussian homogeneity regions with fluctuating
size. A log-normal distribution is a minimal positive-definite model for
multiplicative source-scale fluctuations, for example from multiplicity,
geometry, flow, and \(m_T\) selection effects.

\subsection*{5. Normalized Gamma-Mixture / q-Gaussian}

The effective q-Gaussian form is
\begin{equation}
  S_{\rm raw}(x)
  =
  \left(
  1+\frac{x^2}{\nu R^2}
  \right)^{-(\nu+1)/2}.
\end{equation}

For a radial source distribution:
\begin{equation}
  S_{\rm raw}(r)
  =
  r^2
  \left(
  1+\frac{r^2}{\nu R^2}
  \right)^{-(\nu+1)/2}.
\end{equation}

\begin{center}
\begin{tabular}{lll}
\toprule
Parameter & Meaning & Comment \\
\midrule
\(R\) & Characteristic source scale & Sets the central width \\
\(\nu\) & Tail / fluctuation parameter & Smaller \(\nu\) gives heavier tails; large \(\nu\) approaches Gaussian \\
\bottomrule
\end{tabular}
\end{center}

This form can be interpreted as an effective result of Gaussian scale
fluctuations, for example when a positive-definite width or inverse-width
parameter follows a gamma-like distribution. It is a compact two-parameter
description of a sharp core plus long tail.

\subsection*{Relation to the Gamma Special Function}

The word ``gamma'' in ``gamma-mixture'' does not mean that the fitted source
is the gamma special function itself. The connection is indirect:

\begin{equation}
  \Gamma(k)
  =
  \int_0^\infty t^{k-1}e^{-t}\,dt
\end{equation}

is the gamma special function. It appears as the normalization constant of the
Gamma probability distribution:
\begin{equation}
  P(\beta;k,\theta)
  =
  \frac{1}{\Gamma(k)\theta^k}
  \beta^{k-1}e^{-\beta/\theta},
  \qquad \beta > 0.
\end{equation}

In the source-shape model, the gamma-mixture idea is that a positive-definite
Gaussian scale parameter fluctuates. For example, define an inverse width
parameter
\begin{equation}
  \beta \sim \frac{1}{R^2}.
\end{equation}

If \(\beta\) follows a Gamma distribution, then the observed effective source
shape is obtained by averaging Gaussian shapes over this fluctuation:
\begin{equation}
  S(x)
  =
  \int_0^\infty d\beta\,
  P(\beta;k,\theta)
  \exp\!\left(-\frac{\beta x^2}{2}\right).
\end{equation}

This integral gives a heavy-tailed q-Gaussian / Student-t-like form:
\begin{equation}
  S(x)
  \propto
  \left(
  1+\frac{x^2}{\nu R^2}
  \right)^{-(\nu+1)/2}.
\end{equation}

Thus the logical chain is
\begin{equation}
  \Gamma(k)
  \rightarrow
  \text{Gamma distribution}
  \rightarrow
  \text{Gaussian scale mixture}
  \rightarrow
  \text{q-Gaussian / Student-t-like source shape}.
\end{equation}

In other words, the fit function is not the gamma special function
\(\Gamma(x)\). Rather, it is the effective source shape obtained when the
Gaussian source scale, or inverse scale, fluctuates according to a
Gamma-distributed positive variable.

\section*{Parameter Counting}

\begin{center}
\begin{tabular}{lc}
\toprule
Model & Number of free shape parameters \\
\midrule
Normalized double Gaussian & 3 \\
Normalized symmetric Levy & 2 \\
Normalized Gaussian + exponential tail & 3 \\
Normalized log-normal radius mixture & 2 \\
Normalized gamma-mixture / q-Gaussian & 2 \\
\bottomrule
\end{tabular}
\end{center}

After normalization, none of these models contains a free overall amplitude or
pedestal. Therefore, improvements in \(\chi^2/{\rm NDF}\) are more directly
related to the shape flexibility and physical structure of the model rather
than to arbitrary normalization freedom.

\section*{Interpretation Strategy}

The log-normal and gamma-mixture models should not be interpreted as adding
new pion sources. Instead, they model the observed primordial pair source as a
continuous average over source scales:
\begin{equation}
  S_{\rm obs}(x)
  =
  \int dR\,P(R)\,G(x;R).
\end{equation}

This is consistent with centrality mixing, event-by-event geometry
fluctuations, multiplicity fluctuations, \(m_T\)-selection effects, and flow
geometry variations. The non-Gaussian tail is therefore treated as a property
of the averaged homogeneity region, not as a separate physical emission
component.

\end{document}
